\section{Methodology}

This section walks through the full pipeline of our MSME viability assessment system. We cover the data we used, how we built the five-class label, which models we trained and why, how we added SHAP explanations, the counterfactual recommendation logic, and finally how everything was put together into a working deployed system.

\subsection{Dataset}

For this study we used the U.S.\ Small Business Administration (SBA) National dataset compiled by Li et al.\ \cite{8}. It has 899,164 loan records, all originated somewhere between 1987 and 2014. Each record originally came with 27 variables --- things like loan terms, how much was disbursed, borrower details, and most importantly the final loan outcome, which is stored in a field called \texttt{MIS\_Status} and is either ``Paid in Full'' or ``Charged Off.''

We did not use all 27 variables. After going through the data and thinking about what actually matters for viability (rather than just what was available), we narrowed it down to 11 features. These fell naturally into four groups.

The first group covers the loan itself. \textit{Term} tells us how long the loan is in months and can range anywhere from 0 to 480. \textit{DisbursementGross} is the total dollar amount that was actually handed over. \textit{SBA\_Appv} is the portion that the SBA guaranteed, and \textit{GrAppv} is the gross amount that was approved for the loan overall. All four of these are continuous numerical features.

The second group describes the enterprise. \textit{NoEmp} is the number of employees, which can go from 0 up to 10,000. \textit{NewExist} is a simple flag where 1 means the business already existed and 2 means it was newly established. \textit{UrbanRural} identifies whether the business sits in an urban area (1), rural area (2), or is undefined (0).

Two features capture employment impact. \textit{CreateJob} is the number of new jobs the business planned to create and \textit{RetainedJob} is the number of existing positions it was keeping. Both are integer counts.

The last two features, \textit{RevLineCr} and \textit{LowDoc}, are binary flags related to the loan program. The first indicates whether the loan involved a revolving line of credit and the second flags whether it went through the low-documentation program.

We split the data 80/20 in a stratified fashion, which gave us 720,779 records for training and 178,385 for testing. Continuous features were standardized with \texttt{StandardScaler} from scikit-learn. One thing we were careful about here --- the scaler was fitted only on training data and then applied to both splits, not fitted on the full dataset. That prevents the test set from leaking information back into preprocessing.

Because financial datasets almost always have imbalanced classes (far more loans get paid off than charged off), we applied SMOTE \cite{12} on the training set. This creates synthetic examples of the smaller classes so the models do not just learn to predict the majority class and call it a day.

\subsection{Five-Class Health Taxonomy}

This is probably the most important methodological decision in the paper. Instead of treating loan outcome as a binary variable (which is what every other study does), we constructed a composite health score that captures more nuance about each enterprise.

The idea is straightforward. A business that defaulted on a tiny 12-month loan with zero employees is clearly in worse shape than a business that also defaulted but had a 10-year term, created 30 jobs, and was approved for \$400,000. Binary classification treats both of these as the same ``failure'' --- our taxonomy does not.

We compute a health score $H$ for every record. $H$ is a number between 0 and 100 and it is the sum of four components:

\begin{equation}
H = S_{\text{outcome}} + S_{\text{term}} + S_{\text{job}} + S_{\text{loan}}
\end{equation}

Here is what each one captures:

$S_{\text{outcome}}$ carries 40 out of the 100 points. If the loan was paid in full, the record gets all 40. If it was charged off, it gets 0. We gave this component the largest weight because at the end of the day, whether the business actually paid back its loan is the strongest signal we have about viability. We did experiment briefly with giving this component 30 and 50 points instead, but 40 produced the most balanced distribution across the five resulting classes.

$S_{\text{term}}$ is worth 20 points and comes from the loan duration. We capped it at 240 months (which is 20 years) and then scaled linearly: $S_{\text{term}} = \min(\text{Term}, 240) / 240 \times 20$. The reasoning is that longer terms suggest the lender believed the business had staying power.

$S_{\text{job}}$ is also 20 points and combines \textit{CreateJob} and \textit{RetainedJob}, capped at 50 total jobs: $S_{\text{job}} = \min(\text{CreateJob} + \text{RetainedJob}, 50) / 50 \times 20$. Businesses that create and retain more employment are generally more productive and more viable.

$S_{\text{loan}}$ takes the gross approved amount and scales it the same way, capped at \$500,000: $S_{\text{loan}} = \min(\text{GrAppv}, 500000) / 500000 \times 20$. A larger approved amount signals that the lending institution had confidence in the enterprise.

Figure~\ref{fig:health_score} provides a visual summary of how the four component scores combine into the composite health score and how the resulting value maps to the five viability classes.

\begin{figure}[htbp]
\centering
\includegraphics[width=0.95\linewidth]{figures/health_score_composition.png}
\caption{Composite health score construction. Four domain-informed component scores ($S_{\text{outcome}}$, $S_{\text{term}}$, $S_{\text{job}}$, $S_{\text{loan}}$) are summed to produce $H \in [0, 100]$, which is then discretized into five ordinal viability classes.}
\label{fig:health_score}
\end{figure}

Once every record has its health score, we bin them into five classes as shown in Table~\ref{tab:health_taxonomy}.

\begin{table}[h]
\centering
\caption{Five-class health taxonomy.}
\label{tab:health_taxonomy}
\begin{tabular}{lll}
\hline
\textbf{Health Score Range} & \textbf{Class Label} & \textbf{Interpretation} \\
\hline
$[0, 25)$    & Critical (0) & Defaulted, short term, low employment, small loan \\
$[25, 40)$   & At-Risk (1)  & Defaulted but with some positive indicators \\
$[40, 60)$   & Stable (2)   & Performing loan with moderate features \\
$[60, 75)$   & Growing (3)  & Performing loan with strong employment and size \\
$[75, 100]$  & Thriving (4) & High-performing on all dimensions \\
\hline
\end{tabular}
\end{table}

The boundaries use half-open intervals on the left. So a score of exactly 25 goes into At-Risk, not Critical.

Why does this matter? Because now a loan officer looking at the system output can tell the difference between a business that needs immediate intervention and one that just needs some moderate restructuring. Binary classification cannot make that distinction. And this graduated labeling is what makes the counterfactual recommendation engine in Section~3.5 possible --- you need multiple classes to ask the question ``what would it take to move up one level?''

\subsection{Model Architectures}

We went with two gradient boosting frameworks for our primary classifiers --- XGBoost \cite{10} and LightGBM \cite{11}. We did train a Random Forest \cite{9} too, mostly to have something to compare against, but it never made it into the final deployed system. The saved model files were just too large and took noticeably longer to load during serving, which was a dealbreaker for us.

XGBoost and LightGBM are both gradient boosting methods, but comparing them side by side reveals some interesting architectural differences. Starting with XGBoost --- the way it works is that it grows one decision tree, looks at where that tree made errors, and then grows a second tree whose entire job is to fix those errors. Then a third tree fixes what the first two still get wrong, and so on. The part that really sets it apart from earlier boosting methods is how it handles overfitting. Its loss function has an explicit regularization term baked in:

\begin{equation}
\text{Obj} = \sum_{i} L(y_i, \hat{y}_i) + \sum_{k} \Omega(f_k)
\end{equation}

The $L$ is your standard multi-class cross-entropy. But then $\Omega(f_k) = \gamma T + \tfrac{1}{2}\lambda \|w\|^2$ does double duty --- the $\gamma T$ part discourages trees from having too many leaves (keeps them shallow), and the $\lambda \|w\|^2$ part keeps the leaf weight values from blowing up. There is also the fact that XGBoost takes a second-order Taylor approximation of the loss to figure out the best leaf weights analytically, rather than relying on simple gradient descent. Column subsampling at each split and a shrinkage parameter (the learning rate) give you more ways to avoid overfitting.

LightGBM came along later and was designed to handle bigger datasets without the training time penalty. It introduced two innovations that we found quite impactful during our work. GOSS, or Gradient-based One-Side Sampling, is the first one --- the basic insight is that not all training samples matter equally. Samples with large gradients are the ones the model is struggling with, so GOSS keeps all of those while only taking a random subset of the easy samples. It sounds simple, but it cuts training time by a significant margin. The second trick is called EFB, short for Exclusive Feature Bundling. The idea behind it is fairly intuitive --- in many datasets, you have features that are almost never active at the same time. Think of two binary flags where if one is 1 the other is almost always 0. EFB groups those kinds of features together, and by doing so, the number of features the tree-building process needs to loop through goes down considerably. According to the original LightGBM paper, putting GOSS and EFB together can speed up training by as much as 20 times on standard benchmarks, and the accuracy stays essentially unchanged \cite{11}.

After training, we serialized both models using \texttt{joblib} along with the \texttt{StandardScaler} that had been fitted on the training data. Packaging the scaler alongside the model was a small but important detail --- it ensures that at inference time, incoming data goes through the exact same transformation pipeline that was used during training. We learned early in development that mismatched scaling between training and serving is a surprisingly easy mistake to make and can silently degrade predictions.

\subsection{SHAP-Based Explainability}

From the beginning we were committed to the idea that no prediction should leave the system without an explanation attached. We did not want a black box that just says ``Critical'' with nothing else. A loan officer needs to know what specifically about this application drove the model to that conclusion.

SHAP \cite{17} was our tool of choice for this, and more specifically TreeSHAP \cite{18} --- a variant that was purpose-built for tree ensemble models. Why TreeSHAP and not something like LIME? Two reasons. The first is stability --- LIME generates explanations by perturbing inputs randomly, which means running it twice on the same input can give slightly different results. That is a problem when you need consistent, auditable explanations. TreeSHAP produces exact Shapley values and will always return the same answer for the same input. Speed was the other deciding factor. If you try to compute Shapley values the textbook way, you end up iterating over every possible subset of features --- and with 11 features, that blows up fast. TreeSHAP sidesteps this entirely by exploiting how decision trees are structured internally, which brings the computation down to polynomial time. In our case that meant we could generate explanations within a few milliseconds, which was essential for keeping the API responsive.

For the record, Shapley values come with some nice guarantees from game theory. They satisfy local accuracy, which means the individual feature contributions actually sum up to the model's prediction. They satisfy missingness, so features that the model never sees get attributed zero importance. And they satisfy consistency --- if you retrain the model and a feature becomes more influential, its Shapley value goes up, never down. These are not just academic properties; they matter in practice because they mean the explanations will not behave in counterintuitive ways.

How we actually implemented it: the \texttt{shap.TreeExplainer} object is built once, right when the FastAPI application starts, and then kept around as a singleton for the lifetime of the server process. Creating this explainer is not cheap --- it needs to walk through the entire tree structure and precompute internal data --- so we wanted to make sure it only happens once. During our early prototyping, we had it reinitializing on every request, and the response time was noticeably worse. Caching it as a singleton brought the explanation latency down to something acceptable for production use.

The flow for each request goes like this. The 11 raw feature values come in and get passed through the cached \texttt{StandardScaler}. Then the scaled features go into the XGBoost model, and out comes a predicted class. We then hand that same scaled input to the TreeExplainer and ask for the SHAP values specific to whichever class was predicted. What we get back is a vector of 11 numbers --- one per feature --- where positive values mean that feature pushed the model toward this particular class, and negative values mean it pushed away. We then sort by absolute magnitude and pull out the top three contributors on each side. So the final explanation contains the three biggest factors supporting the classification and the three biggest factors working against it.

To give a concrete example of what the user actually sees: instead of just ``Status: Critical,'' the output might read something along the lines of ``the short loan term was the primary driver of this assessment (SHAP contribution of $+0.58$), while the number of retained jobs worked in the opposite direction, pulling the score away from Critical (SHAP of $-0.05$).'' That level of detail gives practitioners something they can actually reason about and discuss with the applicant, and it provides auditors with a paper trail they can review.

\subsection{Counterfactual Recommendation Engine}

There is a question that SHAP, by itself, cannot answer. Knowing that a short loan term hurt the score is useful, but a loan officer sitting across from an applicant wants to know: okay, so by how many months should we extend it? That is where the counterfactual engine comes in. We built it because we wanted the system to go past diagnosis and into actual guidance --- something we could not find in any existing MSME assessment tool. Two bodies of work shaped our thinking here: Mothilal et al.\ \cite{20} and their DiCE framework for generating diverse counterfactual instances, and the survey by Karimi et al.\ \cite{21} on algorithmic recourse more broadly.

That said, we did not implement DiCE directly. DiCE uses gradient-based or random optimization to search for counterfactuals, and while that can find more optimal solutions in theory, it has a problem for our particular use case: it is stochastic. Run it twice on the same input and you might get a different suggestion. In financial decision-making, that is not acceptable. If two different loan officers assess the same application on the same day, the system absolutely must give them the same recommendation. So we went with a deterministic grid-search approach instead. It sacrifices some theoretical optimality for guaranteed reproducibility, and for our context that tradeoff was easy to make.

Before getting into how the algorithm works, there is an important constraint to address --- not every feature is something that can actually be changed in real life. We went through all 11 features and sorted them into two buckets. Six features made it into the mutable category: \textit{Term} (because you can renegotiate loan duration), \textit{DisbursementGross}, \textit{SBA\_Appv}, and \textit{GrAppv} (all three are dollar amounts that can be revised during restructuring), and \textit{CreateJob} and \textit{RetainedJob} (employment commitments that a business owner can plan around). The remaining five --- \textit{NoEmp}, \textit{NewExist}, \textit{UrbanRural}, \textit{RevLineCr}, and \textit{LowDoc} --- we locked as immutable. You cannot retroactively turn a new business into an existing one, or move it from a rural area to an urban one, so generating recommendations involving those changes would be misleading.

The search process has two phases. To make it concrete, imagine a business has been classified as Critical and the question is: what is the smallest change that would push it into At-Risk?

\textit{Phase~1} looks at single-feature changes. We iterate over each of the six mutable features and, for each one, try a predefined list of values. For \textit{Term} (and other integer features), those are additive changes: add 12 months, add 24, add 36, going up to 240. For dollar-amount features, we use multipliers: scale the original by $0.9\times$, $1.1\times$, $1.2\times$, up to $2.0\times$. Every time we try a perturbation, we scale the modified input with the \texttt{StandardScaler}, feed it through the model, and check if the predicted class changed. When several perturbations do the job, we pick whichever one changed the feature by the smallest absolute amount.

\textit{Phase~2} only runs if Phase~1 fails. It tries pairs of features, but to keep computation from exploding, we limit it to the first five perturbation values in each feature's grid. So instead of testing every possible pair across every possible value, we test all pair combinations but only with a reduced grid. In practice this still covers the range of realistic interventions and runs in well under 10 milliseconds. We chose five as the cutoff after running some experiments during development --- going beyond five did not yield additional feasible counterfactuals in any of our test runs but did increase computation time meaningfully.

Figure~\ref{fig:cf_flowchart} summarizes the full counterfactual generation process from input to recommendation output.

\begin{figure}[htbp]
\centering
\includegraphics[width=0.85\linewidth]{figures/counterfactual_flowchart.png}
\caption{Counterfactual recommendation algorithm. Given an input feature vector and predicted class $c$, the engine first attempts single-feature perturbations (Phase~1). If no single change achieves the target class $t = c + 1$, pairwise perturbations are explored (Phase~2). The minimum-displacement feasible counterfactual is returned.}
\label{fig:cf_flowchart}
\end{figure}

What comes out at the end is a structured recommendation. For each feature the engine decided to change, the output includes four pieces of information: what the original value was, what the new recommended value is, the human-readable name of the feature (so it says ``Loan Duration'' rather than ``Term''), and whether it was increased or decreased. If the grid search exhausted all possibilities without finding anything that works, the engine returns an explicit message saying no feasible recommendation was found. It never fabricates a suggestion when none exists.

We want to highlight three things about why this design suits financial applications specifically. Reproducibility was already mentioned --- same input, same output, every time. Speed is the second one. Because there is no iterative optimization loop and the perturbation grid has a fixed size, the engine is fast enough to run in real-time during a client interaction. And the third thing is simplicity of the output. The recommendations involve at most two feature changes. A loan officer reading the suggestion does not have to mentally juggle five simultaneous modifications --- it is one or two concrete changes, which makes it much easier to discuss with the applicant and to actually follow through on.

\subsection{System Architecture}

Deciding how to deploy everything was its own challenge. Early on we considered building a monolithic application where the ML model, the SHAP explainer, the counterfactual engine, and the user interface all lived in one codebase. But two papers made us reconsider that idea: Sculley et al.\ \cite{26} write extensively about the hidden technical debt that accumulates when ML components get tangled with application code, and Paleyes et al.\ \cite{27} catalog deployment failure modes that frequently stem from exactly this kind of tight coupling. So we went with a clean two-tier split.

The backend runs on FastAPI. We chose FastAPI over Flask because of its built-in support for async handling and because Pydantic validation let us define strict schemas for every incoming request without writing manual checks. There are five endpoints in total. \texttt{POST /predict} takes a single set of 11 features and returns the classification along with the probability distribution over all five classes. \texttt{POST /predict/batch} accepts a CSV upload for portfolio analysis --- useful when someone wants to run the model across hundreds of applications at once. \texttt{POST /explain} returns the SHAP breakdown. \texttt{POST /recommend} generates the counterfactual suggestion. And \texttt{GET /health} is there for monitoring, so deployment tools can check whether the server is up. All five endpoints sit behind API key authentication through the \texttt{X-API-Key} header. On top of auth, Pydantic handles input validation for us. We defined strict ranges for every field --- for example, if someone sends a Term value of 600 or a negative number for NoEmp, the request gets bounced back with a descriptive error before the model ever sees it. This was worth setting up because without it, garbage inputs could produce predictions that look plausible but are meaningless.

We also built in persistent logging from fairly early on. Originally it was not part of the plan, but we quickly realized that in a financial application, you absolutely need to be able to look back at what the system told someone three months ago. So every prediction --- along with the full input vector, the model name, the class probabilities, and whatever recommendation was generated --- goes into a SQLite database through SQLAlchemy. It is nothing fancy from a database perspective, but it gives us a complete paper trail. Down the line, the same logs can be used to catch model drift if the distribution of predictions starts shifting unexpectedly. We also added a \texttt{GET /analytics} endpoint on top of the database so the Streamlit dashboard can pull aggregate statistics without needing direct database access.

For the user interface we went with Streamlit, mainly because it lets you build data-focused dashboards quickly without writing frontend JavaScript. But the important architectural decision here is that the Streamlit app does not run any ML code. It has no idea how XGBoost or SHAP works internally --- it just knows how to call the FastAPI endpoints and display whatever JSON comes back. There are three views a user can access. The first is the single-assessment page, where you punch in the 11 features and see the prediction alongside a SHAP bar chart and the recommendation. The second is a batch mode that takes a CSV file and processes the whole thing at once. The third pulls historical data from the analytics endpoint. One thing we are glad about with this separation is that the frontend could be rebuilt from scratch (say, as a React web app or even a mobile interface) and the backend would not need a single change. That kind of independence has already saved us time when we were tweaking the dashboard layout --- none of those changes required restarting the backend server.
